% =====================================================================
%  proposal_original.tex   —   FROZEN PRISTINE TEMPLATE  (DO NOT EDIT)
%
%  This is the untouched UA project-proposal / SOW template, with all of
%  the original "red italics" guidance intact. It exists so you can always
%  toggle back to the pristine reference on the site. The CI compiles it to
%  proposal_original.pdf, which never changes because this file never changes.
%
%  ==> Edit your real proposal in  proposal.tex,  NOT here. <==
% =====================================================================
\documentclass[11pt]{article}
\usepackage{forge-proposal}
\renewcommand{\showguidance}{1}   % guidance ON for the reference copy

\begin{document}

% ---- Title block (mirrors the Word template's title page) -----------
\begin{center}
  {\Huge\rmfamily\bfseries Your Project Name}\\[6pt]
  {\Large\rmfamily\color{forgemute} Project proposal \& Statement of Work}\\[22pt]
  {\large PM's name, Project Manager}\\[12pt]
  {\large List remaining team members here, one name per line,\\
   first name, last name, alphabetically by last name}\\[12pt]
  {\large POTENTIAL ADVISORS: \placeholder{LIST POTENTIAL ADVISORS FOR YOUR GROUP HERE}}\\[12pt]
  {\large Date: Month, Day, Year}
\end{center}

\vspace{10pt}

\begin{guidance}
Use the statements in red italics to help you write this document. \textbf{Delete all the red italics when you are done.}

This document is both a project proposal and statement of work document. It states, ``Here's the problem we have identified and this is our plan to develop a solution.'' It also clearly details what you will require from the mentor to be successful. It is the contract between you and the mentor which summarizes who will do what and the requirements for success.

SOW: The Product Specification defines your product in detail. What you are making, how it works, how it is used and how it will be tested.

Each section shall have a named author, but all team members are responsible for the overall document. Meaning one team member writes the section and the other team members review it for errors. List the author of each section immediately following the heading and in the author table in the grading rubric at the end. Example: ``This section was written by John Doe.'' It is the team and each member's responsibility to divide the document equitably. Individual scores will be determined by the quality and quantity of each person's work.

Remember the goal is to convey all the information --- ``not applicable because\ldots'' is relevant information. Use bulleted lists and white space smartly. Include figures and tables as required and necessary. For each table and figure, write a brief introductory sentence (``As shown in Table~2 five different pages are used\ldots''). Put the detail in the figures and tables, keep the paragraphs short and do not duplicate information needlessly. Your goal is to convey the full amount of information in the shortest possible document.

You need to maintain a Revision History Table like the one below (not this exact table --- change the dates and entries to match your project).
\end{guidance}

\clearpage

% ---- Revision History -----------------------------------------------
\section*{Revision History}
\renewcommand{\arraystretch}{1.25}
\begin{tabularx}{\linewidth}{|c|Y|c|}
\hline
\textbf{Version} & \textbf{Summary of Changes} & \textbf{Date} \\ \hline
\textit{0.1} & \textit{Section owners assigned; draft Introduction added} & \textit{mm/dd/yy} \\ \hline
\textit{0.5} & \textit{Template info removed; first draft of all sections completed} & \textit{mm/dd/yy} \\ \hline
\textit{0.9} & \textit{All sections updated during team review} & \textit{mm/dd/yy} \\ \hline
\textit{1.0} & \textit{Version submitted for signatures} & \textit{mm/dd/yy} \\ \hline
\end{tabularx}

\begin{guidance}
The revision history table documents major revisions to the document. In addition to showing the evolution of the project it also provides a way for reviewers to quickly see which sections have changed so they do not have to reread the entire document each time. It is a high-level summary --- minor edits, spelling corrections, etc.\ should not need to be summarized. It is your responsibility to determine which sections have significant changes that merit review and decide when it makes sense to declare a revision number. Less detail is required for early revisions. For a shorter document such as the SOW you may only have 2 or 3 versions.
\end{guidance}

% ---- Table of Contents ----------------------------------------------
\begingroup
\renewcommand{\contentsname}{Contents}
\tableofcontents
\endgroup

\begin{guidance}
The table of contents above is generated automatically. Use a page break before the Executive Summary, then start a new page where it makes sense for clarity, balanced with keeping excessive white space out of your document.
\end{guidance}

\clearpage

% =====================================================================
%  1. EXECUTIVE SUMMARY
% =====================================================================
\propsec{1}{Executive Summary}
\sectionauthor{(complete this with the author of each section)}

\begin{guidance}
The executive summary is a short review of the important aspects of the project. It is very similar to the elevator pitch, but slightly longer with more detail. This section should be about 3 paragraphs and about three-quarters of a page.

1st paragraph --- Introduce your product. What is your product and its key features? (``Our product [is a high-efficiency controller / web-based application / wireless smart sensor / \ldots] that [displays / computes / transforms / \ldots] [temperature data / DC current / \ldots] into [easy-to-read graphs / control signals / wireless data] while [lowering cost / improving response time / achieving high efficiency / \ldots].'')

2nd paragraph --- Explain why the product is needed and what problem it is solving. If current solutions exist, explain how your solution is better and unique. Establish ``Why?'' the product will benefit the user. Find credible sources and cite them in a footnote. Don't say ``everyone knows flooding is an issue\ldots''; research and say ``according to [credible reference] flooding in 2019 caused \$32M in damages in Acme County; our product will aid in the detection\ldots''

3rd paragraph --- Describe the process of development: what work will you be doing, who is involved, where the project work will take place, and what you will have completed at the end of the semester.

Include a table as shown. Introduce the table and point out anything significant, e.g.\ ``Table~1 shows the preliminary division of responsibilities among the team members.''

Executive summary max 300 words.
\end{guidance}

\begin{table}[H]
\centering\renewcommand{\arraystretch}{1.25}
\begin{tabularx}{\linewidth}{|l|Y|}
\hline
\textbf{Team Member} & \textbf{Feature responsibility} \\ \hline
\textit{Violet Parr} & \textit{User Interface: Touchscreen and Display} \\ \hline
\textit{Bob Parr} & \textit{Wireless Communications} \\ \hline
\textit{Dashiell Parr} & \textit{Ambient Sensors: Temperature and Angle} \\ \hline
 & \textit{Power: Battery, Charging, Power Budget} \\ \hline
\textit{Helen Parr} & \textit{Remote Management Application (Web based)} \\ \hline
\end{tabularx}
\caption{Preliminary Subsystem Responsibilities}
\end{table}

% =====================================================================
%  2. USER / MARKET RESEARCH
% =====================================================================
\propsec{2}{User/Market research}
\begin{guidance}
This section describes who your product is for and what they want. One of the worst things product developers can do is create a product that no one wants, needs, or will pay for. Look at the macro market to evaluate existing products and market opportunity, and the micro market by performing empathy interviews with potential users.

Overall market: Do some basic research about the size and scope of the market your product will be in.

Existing competitors: List your closest competitors and how your product will be differentiated from them.

User insights: Use empathy interviews to identify needs and pain points, and detail how your product will address them.
\end{guidance}

% ---- 2b ----
\propsec{2b}{Literature Review/Market research}
\begin{guidance}
This section describes who your research project is for, what they want, and what other work has been done in your topic area. It illustrates that you understand the existing research and prior work.

Overall research landscape: Do some basic research about existing research, cite papers, and clarify how your project will contribute new knowledge or use novel approaches on existing data sets.

User insights: Listen to potential users of your research output; list the key insights and pain points discovered and detail how your project will address them.
\end{guidance}

% =====================================================================
%  3. PRODUCT FEATURES
% =====================================================================
\propsec{3}{Product Features}
\begin{guidance}
This section describes the scope and features of the product created by the project --- what the project will PRODUCE.

The key objectives are to:
\begin{itemize}[nosep,leftmargin=1.4em]
  \item list all the key features of the product (``What does it do?'') and who will do them. Include enough work to fulfill the class requirements, but not so much as to create a recipe for failure. Include a minimum feature list and a stretch feature list prioritized by functionality.
  \item summarize the parameters for each feature.
  \item briefly elaborate on the relationship between the product and the business need it addresses. (``Why does the user want this feature? What benefit does it provide?'')
\end{itemize}

A feature is an item that would be listed on the sales pitch to get a customer interested --- a description of what the product accomplishes, not engineering specs. Limit the scope to only what you can accomplish; this also protects against feature creep. Use tables and bulleted lists to keep it concise.
\end{guidance}

\propsub{}{Feature 1: Wireless Transmission of Data}
\sectionauthor{the owner of the subsystem that implements this feature should write this section}
\begin{guidance}
Write a sentence to briefly describe the feature and introduce the table of parameters. Define the requirements --- usually min and max measurable values. Example: ``The system will transmit the water level readings from the remote unit located at the retention pond to the main unit located in the construction office. The requirements for transmission are summarized in Table~2.''
\end{guidance}

\begin{table}[H]
\centering\renewcommand{\arraystretch}{1.25}
\begin{tabularx}{\linewidth}{|Y|c|c|Y|}
\hline
\textit{Parameter} & \textit{Min} & \textit{Max} & \textit{Comments} \\ \hline
\textit{Transmission Distance} & \textit{n/a} & \textit{500\,m} & \textit{Without obstructions, e.g.\ trees} \\ \hline
\textit{Transmission Time} & \textit{n/a} & \textit{200\,ms} & \textit{1 full data packet} \\ \hline
\textit{Wireless Transmitter Power Consumption} & \textit{1\,mW} & \textit{500\,mW} & \textit{Min during sleep mode, max during active transmission} \\ \hline
\end{tabularx}
\caption{Wireless Transmission Parameters}
\end{table}

\propsub{}{Feature 2: Temperature Sensing and Logging}
\sectionauthor{the owner of the subsystem that implements this feature should write this section}
\begin{guidance}
Repeat the format from Feature~1 as needed. You should have a minimum of 3 key features and a maximum of 8. Most projects will have 4 or 5.
\end{guidance}

% ---- 3b ----
\propsec{3b}{Research Project Deliverables}
\begin{guidance}
This section describes the scope and features of the analysis created by the project --- what the project will PRODUCE. Write a draft, but this is what your mentor will especially help you with. This is often the breakdown of what will and will not be considered a passing project, so make this detailed.

\textbf{Final Presentation Format.} How will your final project be submitted? Will it culminate in a paper with introduction, methods, and results sections and at least 5 peer-reviewed citations, under 10 pages, with at least two informative graphics? Will it culminate in an R Shiny application available on GitHub?

\textbf{What Analysis Is Being Run?} Once the data is in, what analyses exactly should be run? E.g.\ what type of neural network, what type of correlation analysis, etc.

\textbf{What Accuracy Is Expected?} If you're running a model, what do you expect --- 100\% accuracy? 65\% across multiple outputs? An $R^2$ of 0.6? What happens if the model isn't very accurate?

\textbf{What if the Analysis doesn't work?} Is a null result acceptable, or do you need to do work to circumvent that?

\textbf{What if the Data Isn't Available?} For projects requiring data scraping --- what happens if your data is not scrapable? How will you salvage your project?
\end{guidance}

% =====================================================================
%  4. PROJECT TIMELINE & GANTT CHART
% =====================================================================
\propsec{4}{Project Timeline \& Gantt Chart}
\begin{guidance}
List the major deliverables and decision points of the project. Track your work over time using the Gantt chart and continually update and submit it in weekly check-ins. Milestones from the course schedule are pre-populated; add the dates for the current semester and your project-specific objectives and deadlines.

A very common mistake is to build the schedule to have a functional project by Demo day, leaving no time for testing and modification (allow 6--8 weeks for testing minimum) and no time to recover if early dates slip. At this stage you should have about 10--15 items in addition to the milestones from the master schedule.
\end{guidance}

\begin{longtable}{|p{10cm}|c|}
\hline
\textbf{Milestone} & \textbf{Date} \\ \hline
\endfirsthead
\hline
\textbf{Milestone} & \textbf{Date} \\ \hline
\endhead
Team Formation & 1/29/24 \\ \hline
 & \\ \hline
 & \\ \hline
Signed proposal & 2/19/24 \\ \hline
 & \\ \hline
 & \\ \hline
 & \\ \hline
 & \\ \hline
 & \\ \hline
 & \\ \hline
Poster Demo & 4/23/24 \\ \hline
 & \\ \hline
iShowcase & 5/1/24 \\ \hline
\caption{Milestone Schedule}
\end{longtable}

\begin{guidance}
Some items to consider as examples: determining sub-system partitioning and assignments; preliminary hardware selection; selecting a microcontroller / programming language / application framework; individual sub-system builds complete; individual sub-system testing complete; complete system build integrating the sub-systems; full-function testing complete; full feature demonstration.

A Gantt chart is a visual representation of the project timeline linked to a schedule with task ownership and dependencies. Refer to the Gantt Chart module on D2L for instructions and resources.
\end{guidance}

% =====================================================================
%  5. ETHICS
% =====================================================================
\propsec{5}{Ethics}
\begin{guidance}
Data ethics is a key consideration in any product we create. Include a completed ethics chart (see examples on D2L). For each question, answer Y / N / M (Yes / No / Maybe) for general use and for the data-breach scenario.
\end{guidance}

\renewcommand{\arraystretch}{1.2}
\begin{longtable}{|c|p{9.5cm}|c|c|}
\hline
\textbf{\#} & \textbf{Question} & \textbf{Generally} & \textbf{Data Breach} \\ \hline
\endfirsthead
\hline
\textbf{\#} & \textbf{Question} & \textbf{Generally} & \textbf{Data Breach} \\ \hline
\endhead
1 & Could a user sell drugs or other illegal items on your platform? & Y/N/M & Y/N/M \\ \hline
2 & Could a user of your platform engage in sex trafficking? & Y/N/M & Y/N/M \\ \hline
3 & Could a user sell class notes or cheat on their homework on your platform? & Y/N/M & Y/N/M \\ \hline
4 & Could a stalker use your project to find someone? & Y/N/M & Y/N/M \\ \hline
5 & Could your app be used to spy on or track individuals? & Y/N/M & Y/N/M \\ \hline
6 & Could your app/software access the camera or microphone and record things without users being aware? & Y/N/M & Y/N/M \\ \hline
7 & If someone uses your platform, could they be re-traumatized or have their mental health impacted in some way? & Y/N/M & Y/N/M \\ \hline
8 & Could your algorithm promote material that would traumatize or upset individuals? & Y/N/M & Y/N/M \\ \hline
9 & Would your users be upset if the data you collect was given to someone else? & Y/N/M & Y/N/M \\ \hline
10 & Could a data leak potentially lead to identity theft? & Y/N/M & Y/N/M \\ \hline
11 & If your site was hacked, would users potentially lose their job, spouse, or family? & Y/N/M & Y/N/M \\ \hline
12 & Should there be an age limitation on your product? & Y/N/M & Y/N/M \\ \hline
13 & Could someone use your product to find, contact, and potentially commit elder abuse? & Y/N/M & Y/N/M \\ \hline
14 & If the data on your platform was breached, could it be used to blackmail the users? & Y/N/M & Y/N/M \\ \hline
15 & Does the existence of your project imply that a particular racial group, gender, religion, or other protected category is inherently bad, gross, or unwanted? & Y/N/M & Y/N/M \\ \hline
16 & Could your product be used to commit hate crimes against a specific group? & Y/N/M & Y/N/M \\ \hline
17 & Does the primary content of your game or algorithm focus on something considered deeply unethical? & Y/N/M & Y/N/M \\ \hline
18 & Does your game or software contain race, gender, or other stereotypes? & Y/N/M & Y/N/M \\ \hline
19 & Could users of your app scam other individuals? & Y/N/M & Y/N/M \\ \hline
20 & Is your particular algorithm biased toward predicting correctly only for one race, gender, or other group? & Y/N/M & Y/N/M \\ \hline
21 & Are the users of your project aware of how their data will be used? & Y/N/M & Y/N/M \\ \hline
22 & What are the possible misinterpretations of your results? & Y/N/M & Y/N/M \\ \hline
23 & Does the use or purchase of your data potentially contribute to a dangerous group or regime? & Y/N/M & Y/N/M \\ \hline
24 & Could your virtual reality environment cause injury to the user? & Y/N/M & Y/N/M \\ \hline
25 & Are your study participants or game players aware that their data will be collected and used? & Y/N/M & Y/N/M \\ \hline
26 & Does your game or app contain addictive design elements without benefit to the user? & Y/N/M & Y/N/M \\ \hline
27 & Does your survey contain an aspect of compulsion or unusually large incentive that would command users to take it even to their detriment? & Y/N/M & Y/N/M \\ \hline
28 & Could your research outcomes harm an individual or entity? & Y/N/M & Y/N/M \\ \hline
\caption{Data Ethics Chart}
\end{longtable}

% =====================================================================
%  6. APPROVALS
% =====================================================================
\propsec{6}{Approvals}
\begin{guidance}
The project proposal needs the approval of everyone to move forward. If approval is denied, rework the proposal and begin the approval cycle again. Consult the faculty advisor often. After the faculty advisor and team are comfortable with the SOW, only then should it go to the sponsor.

Type the names in the ``Approver Name'' column; written signatures only in the ``Signature'' and ``Date'' columns. Give advisors and sponsors at least 3 business days of lead time. Your instructor signs after grading, if the document passes. Keep the language below.
\end{guidance}

\noindent The signatures of the people below indicate an understanding of the purpose and content of this document by those signing it. By signing this document, you indicate that you approve of the proposed project outlined in this Statement of Work, the division of work, the Ground Rules, and that the next steps may be taken to create a Product Specification and proceed with the project.

\smallskip
\noindent This document is based upon and supersedes the \placeholder{<PRD title>} Version~\placeholder{X.X}. Deviations (versus clarifications) from the PRD have been clearly noted. For any requirements not listed in this SOW, the PRD requirements shall remain in effect.

\medskip
\renewcommand{\arraystretch}{1.6}
\begin{tabularx}{\linewidth}{|Y|l|c|c|}
\hline
\textbf{Approver Name} & \textbf{Title} & \textbf{Signature} & \textbf{Date} \\ \hline
\placeholder{Type names here} & Team Project Manager & & \\ \hline
\placeholder{All team members must sign} & Team Member & & \\ \hline
 & Team Member & & \\ \hline
 & Advisor & & \\ \hline
 & Instructor & & \\ \hline
\end{tabularx}

\begin{guidance}
Before you submit to advisor/sponsor for review: refresh the table of contents; verify every table and figure has a number and caption with at least one introductory sentence; clean up white space and page breaks; make sure the revision block is correct.

Before you submit for grading, update the author table below. It need not be 100\% accurate but should give the instructor a good idea of how much each person contributed. For tables and diagrams just put n/a for the word count.
\end{guidance}

\medskip
\renewcommand{\arraystretch}{1.4}
\begin{tabularx}{\linewidth}{|Y|l|c|}
\hline
\textbf{Section} & \textbf{Author} & \textbf{Word Count} \\ \hline
\textit{1. Executive Summary} & \textit{Jane Doe} & \textit{268} \\ \hline
 & & \\ \hline
 & & \\ \hline
 & & \\ \hline
\end{tabularx}

% =====================================================================
%  7. APPENDIX
% =====================================================================
\propsec{7}{Appendix}

\propsub{A.}{Advisor Engagement}

\propsubsub{Project Team Responsibilities}
\begin{itemize}[leftmargin=1.4em]
  \item The Project Manager will set up and facilitate a weekly call/meeting with the Faculty Advisor. The Project Team will provide weekly status updates including upcoming deliverables, critical issues, and any adjustments to the Project Plan.
  \item Documents will be provided to the Faculty Advisor with adequate time for review and signature, with a minimum review time of 3 days prior to the document due date.
  \item Design files will be provided to the Faculty Advisor as requested in an agreed format.
  \item Support requirements will be clearly requested with the dates required and adequate time for fulfillment.
  \item Modification requests to the Project Plan by the Faculty Advisor will be reviewed and agreed to within 1 week of the request.
\end{itemize}

\propsubsub{Faculty Advisor Responsibilities}
\begin{itemize}[leftmargin=1.4em]
  \item The Faculty Advisor will provide knowledge and expertise to help the group stretch their skills.
  \item The Faculty Advisor will participate in a weekly or bi-weekly call/meeting to review status, upcoming deliverables, priorities, issues, and progress to the agreed Project Plan.
  \item The Faculty Advisor will provide document review, feedback, and approval (or rejection, or approval with contingencies) with adequate time for the team to meet course due dates.
  \item The Faculty Advisor will provide feedback on support requirements, including guidance on design decisions, design files, test plans, test procedures, and test results.
  \item The Faculty Advisor shall provide technical advice and guidance, answering inquiries approximately 1 hour per week.
  \item Modifications to the Project Plan by the Project Team will be resolved and documented within 1 week of the request.
  \item Grade the finalized project using a skill-based rubric.
  \item Attend iShowcase in May.
\end{itemize}

\propsub{B.}{Ground Rules}
\begin{guidance}
How the team will conduct the business of completing this project. Each team member must sign the Approvals section indicating their acknowledgement of these Ground Rules. Do not remove items from this list, but you may add items.
\end{guidance}

\noindent As a team and as individual team members, we agree to:
\begin{itemize}[leftmargin=1.4em]
  \item \textbf{Stay focused on our objectives and goals.} Each time the team meets, we will clearly define our objectives and desired outcomes at the beginning of the meeting, and politely remind team members if we are getting off track.
  \item \textbf{``Sidebar'' issues that are relevant but not consistent with the immediate objectives.} Create a ``sidebar'' where off-topic but important matters can be listed and discussed later.
  \item \textbf{Listen when others are speaking.} We will listen and consider others' input before adding our own comments.
  \item \textbf{All viewpoints will have an opportunity to be heard.} We will make an effort to get each team member's viewpoint and ensure no one dominates the discussion.
  \item \textbf{Differences of opinion will be discussed respectfully.} We will identify areas of agreement before assessing disagreement, weigh the merits of different opinions, and agree on a process for choosing a direction. All team members will respect and follow the decision.
  \item \textbf{Look for the good points in new ideas.} We will explore the value in each idea as we assess and select our path forward.
  \item \textbf{Focus on the future, not the past.} We will use past experience to inform decisions but focus the discussion on future objectives; blame is counterproductive.
  \item \textbf{Agree upon specific action items and next steps.} At the end of each meeting we will summarize and agree on specific next steps, action items, and assignments.
  \item \textbf{Accountability.} Each member is responsible for their individual assignments and contribution to the team objectives. We will honor our responsibilities and not let our team members down.
\end{itemize}

\end{document}
